\chapter{Testing and Results}

Everything in this chapter is measured on the held-out 20\% test set --- the same 102 unseen contracts for every model. The numbers were regenerated by a standalone evaluation notebook that loads the trained models from disk, so no state left over from the training session can influence them.

\section{Presence Classification}

\subsection{Model comparison}
Table~\ref{tab:ensemble} and Figure~\ref{fig:ensemble} compare all the presence models on the 4{,}182 (contract, category) test decisions. Two things stand out. First, the windowed transformer \emph{loses on its own} to the full-document TF--IDF baseline (0.694 vs.\ 0.775). Legal vocabulary is distinctive enough that lexical evidence over the whole document beats semantic evidence over windows, and max-pooling inflates false positives exactly the way Section~\ref{sec:mil} predicted. Second, the AND-ensemble --- which reports a clause only when \emph{both} models agree --- filters out each model's uncorrelated false positives and finishes on top (0.776). The margin over the baseline is small on 102 test contracts, so we do not lean on it: Section~\ref{sec:bootstrap} quantifies exactly how small, and finds that this particular gap is indistinguishable from zero. The conclusion we stand behind is methodological: \emph{simple lexical models are strong at document-level presence, and the transformer earns its place as a precision filter, not as a replacement}.

\begin{table}[ht]
\centering
\caption{Presence models on the 20\% test set (micro-F1 over all 4{,}182 test cells).}
\label{tab:ensemble}
\begin{tabular}{lc}
\toprule
\textbf{Model} & \textbf{micro-F1}\\
\midrule
\textbf{Ensemble AND (both must agree)} & \textbf{0.776}\\
TF--IDF baseline (full document)        & 0.775\\
Ensemble AVG (mean probability)         & 0.718\\
Ensemble OR (either fires)              & 0.696\\
Transformer (max-pool) alone            & 0.694\\
\bottomrule
\end{tabular}
\end{table}

\begin{figure}[ht]
\centering
\includegraphics[width=0.75\textwidth]{ensemble_comparison.png}
\caption{micro-F1 of all presence models on the held-out test set.}
\label{fig:ensemble}
\end{figure}

\subsection{How much of that ranking is real?}
\label{sec:bootstrap}
Table~\ref{tab:ensemble} is a leaderboard, and leaderboards invite a reader to treat the ordering itself as a result. On 102 test contracts a margin of 0.776 against 0.775 --- roughly one decision in a thousand --- plainly cannot support that reading, so rather than leave the caveat as prose we measured it.

The natural tool is a \textbf{bootstrap confidence interval}, but the resampling unit matters. The 4{,}182 test decisions are not independent: 41 of them come from the same contract, and a contract that is unusual --- unusually long, unusually drafted --- moves all 41 together. Resampling individual (contract, category) cells would ignore that structure and produce intervals that are far too narrow. We therefore resample \emph{contracts} with replacement (a cluster bootstrap), drawing 102 contracts from the 102 available, recomputing micro-F1 for both models on each resample, and repeating 10{,}000 times. Both models are scored on exactly the same resample every time, so the difference is \emph{paired} --- which is what allows a difference interval much tighter than the two individual intervals would suggest.

\begin{table}[ht]
\centering
\caption{Cluster bootstrap over the 102 test contracts (10{,}000 resamples). The paired difference is what settles each comparison; $P(>0)$ is the fraction of resamples in which the first model wins.}
\label{tab:bootstrap}
\small
\begin{tabular}{lccc}
\toprule
\textbf{Comparison} & \textbf{Difference in micro-F1} & \textbf{95\% CI} & \textbf{$P(>0)$}\\
\midrule
AND-ensemble vs.\ TF--IDF baseline & $+0.0014$ & $[-0.0068,\ +0.0089]$ & 0.626\\
TF--IDF baseline vs.\ transformer  & $+0.0805$ & $[+0.0626,\ +0.0992]$ & 1.000\\
AND-ensemble vs.\ transformer      & $+0.0818$ & $[+0.0654,\ +0.0987]$ & 1.000\\
\bottomrule
\end{tabular}
\end{table}

The three comparisons do not come out the same way, and the difference between them is the point.

\textbf{The ensemble's win over the baseline is not real.} The interval is $[-0.0068, +0.0089]$: it straddles zero, and the ensemble is ahead in only 62.6\% of resamples, barely better than a coin flip. Taken individually the two models score 0.776 (95\% CI $[0.759, 0.792]$) and 0.775 (95\% CI $[0.758, 0.791]$), intervals that almost entirely overlap. On this test set the AND-ensemble and the TF--IDF baseline are \emph{statistically indistinguishable}, and we now state that as a measurement rather than as a hedge. The top two rows of Table~\ref{tab:ensemble} should be read as a tie.

\textbf{The transformer's loss, by contrast, is real.} The baseline beats the windowed transformer by 0.0805, with an interval of $[+0.0626, +0.0992]$ that is nowhere near zero, and it wins in every one of the 10{,}000 resamples. This matters for how the thesis' two empirical claims should be weighted, because they are not on equal footing. The finding we stand behind most firmly --- that a full-document lexical model beats a windowed DistilBERT at document-level presence classification --- is the one that clears the significance test comfortably. The finding we were already cautious about --- that the conjunctive ensemble edges past the baseline --- is the one that does not. That our caution happened to point the right way is reassuring, but it was worth checking rather than assuming, and a reader should not have had to take it on trust.

None of this undercuts the case for the ensemble, because that case was never the 0.001 of micro-F1. The AND rule earns its place through the \emph{error profile} it produces rather than its position on the leaderboard: it converts the transformer's inflated false positives into balanced precision and recall, which is what the next section examines.

\subsection{Confusion-matrix analysis}
Our evaluation protocol called for a full confusion-matrix analysis of the final model. Table~\ref{tab:cm} gives the matrix for the AND-ensemble over all test decisions, Figure~\ref{fig:cm} shows the same matrix as a heatmap, and Table~\ref{tab:metrics} lists every metric derived from it.

\begin{table}[ht]
\centering
\caption{Confusion matrix of the AND-ensemble on the 20\% test set (4{,}182 decisions).}
\label{tab:cm}
\begin{tabular}{lcc}
\toprule
 & \textbf{Predicted Absent} & \textbf{Predicted Present}\\
\midrule
\textbf{Actually Absent}  & TN = 2{,}523 & FP = 311\\
\textbf{Actually Present} & FN = 296     & TP = 1{,}052\\
\bottomrule
\end{tabular}
\end{table}

\begin{table}[ht]
\centering
\caption{Test metrics derived from the confusion matrix.}
\label{tab:metrics}
\begin{tabular}{llc}
\toprule
\textbf{Metric} & \textbf{Formula} & \textbf{Value}\\
\midrule
Accuracy    & $(TP+TN)/\text{all}$          & 85.49\%\\
Precision   & $TP/(TP+FP)$                  & 77.18\%\\
Recall (Sensitivity) & $TP/(TP+FN)$         & 78.04\%\\
Specificity & $TN/(TN+FP)$                  & 89.03\%\\
F1 score    & $2PR/(P+R)$                   & 0.7761\\
\bottomrule
\end{tabular}
\end{table}

\begin{figure}[ht]
\centering
\includegraphics[width=0.6\textwidth]{confusion_matrix.png}
\caption{Confusion-matrix heatmap of the AND-ensemble on the 20\% test set.}
\label{fig:cm}
\end{figure}

Precision and recall come out balanced (77.2\% and 78.0\%), which tells us the AND rule did its job: raw max-pooling gives up precision, and requiring agreement wins it back without sacrificing recall. The specificity of 89.0\% means the system rarely claims a clause that is not there --- in a legal-review assistant, that is the more damaging kind of mistake.

\subsection{Per-category behaviour}
Figure~\ref{fig:percat} shows per-category F1 for the strongest categories. The frequent, formulaic ones are close to solved (\emph{Document Name} 1.000, \emph{Parties} 0.995, \emph{Governing Law} 0.961, \emph{Agreement Date} 0.958). At the other end, the five categories with at most seven test positives (\emph{Most Favored Nation}, \emph{No-Solicit of Customers}, \emph{Price Restrictions}, \emph{Source Code Escrow}, \emph{Unlimited License}) all score 0.000. That number needs context: with two positive examples in the test set, a score of zero means the model missed two contracts, not that it is broken. This tail is a property of the dataset rather than of any architecture, and it is why we quote micro- rather than macro-averaged numbers in our headline results.

\begin{figure}[ht]
\centering
\includegraphics[width=0.85\textwidth]{per_category_f1.png}
\caption{Per-category F1 on the test set for the twelve strongest categories.}
\label{fig:percat}
\end{figure}

\section{Span Extraction}
We evaluated the QA model on all 2{,}667 answer-bearing test windows (Table~\ref{tab:span}). A token-F1 of 0.763, with 82.7\% of predictions overlapping the gold span at F1 $\geq$ 0.5, means that when the clause is in front of the model, it captures most of the right text roughly eight times out of ten. The lower exact-match score did not surprise us: annotated legal spans are long, and their exact boundaries are genuinely debatable (should a leading section number be included?), so content overlap is the metric that carries information here. For scale, the published CUAD baselines at this task use DeBERTa-v3-large, roughly five times the parameters of our distilled backbone.

\begin{table}[ht]
\centering
\caption{Span-extraction results on the 2{,}667 answer-bearing test windows.}
\label{tab:span}
\begin{tabular}{lc}
\toprule
\textbf{Metric} & \textbf{Score}\\
\midrule
Token-level F1              & 0.763\\
Overlap (token-F1 $\geq$ 0.5) & 82.7\%\\
Exact Match                 & 35.5\%\\
\bottomrule
\end{tabular}
\end{table}

\textbf{Scope caveat.} These numbers measure extraction quality \emph{given an answer-bearing window}. The full-document benchmark in the CUAD paper (precision at 80\% recall over all windows, including the answer-free ones) needs a heavier document-stride ranking evaluation that we did not run, so we make no claim on that metric.

\section{Summarization}
On 150 held-out test clauses the fine-tuned FLAN-T5-small scores \textbf{ROUGE-L 0.775}, well above the 0.40 target we set in the proposal. On its own, that number would overstate what the model actually learned. The qualitative comparison in Table~\ref{tab:sumqual} shows what is really going on: the fine-tuned model has learned to \emph{classify} the clause and print its category template word for word. The output is clean, readable, and usually correct --- but it is category-generic, and it ignores the details of the specific clause. Zero-shot FLAN-T5 does the opposite: it engages with the actual clause but only echoes the legalese back without simplifying it. With just 41 unique target sentences, fine-tuning quietly turned summarization into classification. The high ROUGE is real, but what it measures is template reproduction, and we would rather say so than let the metric overclaim. A genuinely abstractive legal summarizer will need a larger, clause-diverse target set and a bigger model.

\begin{table}[ht]
\centering
\caption{Fine-tuned vs.\ zero-shot summarizer on the same test clause (abridged).}
\label{tab:sumqual}
\small
\begin{tabular}{p{3.2cm} p{10.3cm}}
\toprule
\textbf{Input clause} & ``Any attempted assignment or delegation in violation of this section by either party without the prior written consent of the other will be void.''\\
\midrule
Gold template & You cannot transfer this contract to someone else without permission.\\
Fine-tuned    & You cannot transfer this contract to someone else without permission. \emph{(exact template)}\\
Zero-shot     & ``a) Any attempted assignment or delegation in violation of this section \ldots will be void.'' \emph{(echoes the legalese)}\\
\bottomrule
\end{tabular}
\end{table}

\section{Summary Scorecard}
Table~\ref{tab:scorecard} pulls together the held-out results of all three models. Every number was produced by the standalone test notebook running on the saved 80/20-trained models.

\begin{table}[ht]
\centering
\caption{Final held-out test results (408 training / 102 test contracts).}
\label{tab:scorecard}
\small
\begin{tabular}{llll}
\toprule
\textbf{Stage} & \textbf{Model} & \textbf{Metric} & \textbf{Score}\\
\midrule
2A Presence & AND-ensemble (TF--IDF + DistilBERT) & micro-F1  & 0.776\\
2A Presence & AND-ensemble & Accuracy               & 85.49\%\\
2A Presence & AND-ensemble & Precision / Recall     & 77.18\% / 78.04\%\\
2A Presence & AND-ensemble & Specificity            & 89.03\%\\
2B Span     & DistilBERT-QA & token-F1              & 0.763\\
2B Span     & DistilBERT-QA & Exact Match           & 35.5\%\\
1 Summarizer & FLAN-T5-small & ROUGE-L              & 0.775\\
\bottomrule
\end{tabular}
\end{table}

\section{Discussion: the Three Findings}
\begin{enumerate}[itemsep=2pt]
  \item For document-level \emph{presence}, a simple lexical baseline beats a windowed transformer, and a paired cluster bootstrap puts that gap at $+0.081$ micro-F1 with a 95\% interval of $[+0.063, +0.099]$ --- it is not a sampling artefact. The transformer pays for itself only as a precision filter inside a conjunctive ensemble, and even there the ensemble's headline margin over the baseline is statistically a tie; what the ensemble genuinely buys is a balanced error profile, not a higher score.
  \item For \emph{span extraction}, the transformer has no substitute and performs well even in distilled form --- this is where the neural model genuinely earns its computational cost.
  \item For \emph{summarization}, a tiny template seed quietly turns the task into classification. A high automatic metric should not be believed until the outputs have actually been read.
\end{enumerate}
